\chapter{Anexo 2: Aromas de varias sustancias}
En este anexo se lleva un seguimiento de la obtencion de las características de la sustancia que se desea detectar en distintas condiciones. Estos datos se usarán para el entrenamiento de la nariz electrónica y serán clasificadas como ``BUENO''.

\begin{table}[htbp]
	\centering
	\footnotesize
	\setlength{\tabcolsep}{4pt}
	\begin{tabularx}{\textwidth}{@{} >{\raggedright\arraybackslash}p{0.18\textwidth} >{\raggedright\arraybackslash}X >{\raggedright\arraybackslash}X >{\raggedright\arraybackslash}X @{}}
		\toprule
			\textbf{Nombre} & \textbf{Sustancia} & \textbf{Condiciones} & \textbf{Hora inicio; Hora fin; Comentarios} \\
		\midrule
		EXP\_M\_1 & Cacao en polvo (2\,g) & Bote cerrado a temperatura ambiente & {\raggedright\begin{minipage}[t]{\linewidth}\begin{itemize}\itemsep0pt\parskip0pt\parsep0pt\setlength{\topsep}{0pt}
			\item \textbf{Fecha:} 2025--12--11
			\item \textbf{Hora inicio:} 17:16
			\item \textbf{Hora fin:} 17:22
			\item \textbf{Comentarios:}
		\end{itemize}\end{minipage}}\\
		\midrule
		\bottomrule
	\end{tabularx}
	\caption{Registro de mediciones de sustancias y condiciones de muestreo.}\label{tab:registro_sustancias_bueno}
\end{table}